\documentclass[11pt]{article}
\usepackage{../latex/styles/core} 

\bibliography{bibliography}

\title{Hands On DL - MIT 15.776}
\author{Sai Mohan Kilambi}
\date{\today}
\begin{document}

\maketitle

\setcounter{tocdepth}{4}
\tableofcontents

\newpage
\section{Lecture 2 - Training Deep Neural Networks}
Someone asked the question when should you use DL and or something simpler like Binary Trees. The answer was that there are two dimensions
\begin{enumerate}
	\item Is explaining the model more important? Go with simpler models
	\item Is accuracy more important? Then go with that and dont worry how complex the model is. 
\end{enumerate}

By default just choose \textbf{RELU} as your non linear function.

Lets design the layout of the network
\begin{enumerate}
	\item Choose number of hidden layers and number of neurons in each layer - We choose 1 Hidden layer with 16 RELU neurons
	\item Pick the right output layer - Sigmoid when you want output between 0 and 1 like a probability.
\end{enumerate}

If there are no hidden layers its just \textbf{Linear Regression}.

Too many layers can result in overfitting.

If you want to emit a probability in the end then use Sigmoid function. Thats why the output layer is sigmoid. HE was explaining a medical diagnosis problem.

Ok lets assume a network with on hidden layer of 16 nodes. So in this medical diagnosis problem there are 
\begin{enumerate}
	\item 29 inputs
	\item Hidden layer: 16 nodes with weight and bias for each node.
	\item Output layer with 1 weight and 1 bias.
\end{enumerate}

So total number of parameters is $29*16 + 16 + 16*1 +1 = 497$. You see how the parameters are just exploding. 

\begin{lstlisting}[style=pythonstyle]
import keras

#Create the input layer
input = keras.Input(shape=(16,)) 

# Now create the first hidden layer
h = keras.layers.Dense(16, activation="relu")(input)
output = keras.layers.Dense(1, activation="sigmoid")(h)

model = keras.Model(input, output)
\end{lstlisting}

The role of the data is to give you the coefficients.
\end{document}